\chapter{Filtering a sequence of samples}
\MARSChapterHeading{Filtering a sequence of samples}
\label{revision:filtering}

A mixer combines values at the same position in two streams. A filter can combine
several positions in one stream. Imagine replacing each sample by the average of
itself and its seven predecessors. An isolated change then influences eight
outputs instead of one. This changes the signal even when no arithmetic clips.
To implement that operation, we must specify which samples participate, how much
each contributes, and when the calculation loses precision.

The complete \texttt{fir-filter.asm} in \texttt{revision/labs/files} implements an
eight-coefficient finite impulse response filter, abbreviated FIR. Its scalar and
MSA paths compute the same exact sum before rounding once and saturating once.
This chapter assumes the earlier fixed-point and PCM chapters and MSA lane
operations. The independent task uses the same arithmetic but retains only every
second filtered position.

\section{Give each coefficient a time position}

Call the input array $x$ and the eight coefficient codes $c_0,\ldots,c_7$. The
array begins with seven history samples, in chronological order. Its next element
is the first new sample. For output index $i$, the input window is
$x_i,\ldots,x_{i+7}$. The first coefficient multiplies the oldest value in that
window; the last multiplies its newest value. A coefficient applied at one time
position is often called a \emph{tap}.

The function forms the integer numerator
\[
 S_i=\sum_{j=0}^{7}x_{i+j}c_j.
\]
These are signed sixteen-bit input and coefficient codes. A coefficient code $c$
represents the weight $c/32768$, a format commonly called Q15 here. Code 4096
means one eighth, 16384 means one half, and $-32768$ means minus one. Positive
one has no exact signed Q15 code: its scaled integer 32768 exceeds 32767.
The sample remains an integer PCM code; multiplying by a Q15 coefficient gives a
numerator in sample-code units scaled by 32768. Dividing the final sum by that
scale restores sample-code units.

With eight coefficients of 4096, the weights sum to one. The operation averages
the window. A constant input of 16000, including its history, therefore produces
16000. If history is zero and only the first new sample is 16000, the first
window produces 2000. The same impulse remains in each of the next seven windows,
then leaves. Outputs zero through seven are 2000; later outputs are zero.
The response to that single impulse ends after finitely many outputs, giving
this kind of filter its name. Each output depends on a finite input window,
not on a chain of earlier outputs. History values are inputs, not an instruction
to reset the filter at every call.

\begin{figure}[htbp]
  \centering
  \begin{tikzpicture}[alt={Two overlapping eight-sample windows for a single impulse. The first seven samples are zero history. Original sample seven is 16000, followed by zeros. Output zero uses indices zero through seven; output one uses one through eight. The impulse remains in each window through output seven; output eight no longer contains it. With every coefficient one eighth, outputs zero through seven are 2000 and output eight is zero.},x=.67cm,y=.65cm]
    \foreach \j in {0,...,9} {
        \draw (\j,0) rectangle +(1,1);
        \node[font=\scriptsize] at (\j+.5,1.35) {\j};
      }
    \foreach \j in {0,...,6,8,9} {\node[font=\small] at (\j+.5,.5) {0};}
    \node[font=\scriptsize\bfseries] at (7.5,.5) {16000};
    \draw[line width=1pt] (0,-.25)--(0,-.5)--(8,-.5)--(8,-.25);
    \node[anchor=east,font=\small] at (-.15,-.5) {$i=0$};
    \draw[line width=1pt] (1,-.85)--(1,-1.1)--(9,-1.1)--(9,-.85);
    \node[anchor=east,font=\small] at (.85,-1.1) {$i=1$};
    \node[anchor=west,font=\small] at (8.2,-.5) {$2000$};
    \node[anchor=west,font=\small] at (9.2,-1.1) {$2000$};
  \end{tikzpicture}
  \caption{A window moves one sample for each output. Seven history samples make the first complete window available. The impulse initially meets coefficient seven, then coefficient six, and so on until it leaves the window.}
  \label{fig:revision-fir-windows}
\end{figure}


Coefficient order matters when the weights differ. For an impulse of 16384 and
oldest-first codes $1024,2048,3072,4096,5120,6144,7168,8192$, output zero is
4096 because the impulse initially meets the last coefficient. Subsequent outputs
are $3584,3072,2560,2048,1536,1024,512$, then zero. An all-equal averaging filter
cannot reveal a reversed coefficient order; this asymmetric example can.

\section{Keep the exact sum until the chosen conversion}

The required integer result before saturation is
\[
 R_i=\left\lfloor\frac{S_i+16384}{32768}\right\rfloor.
\]
Adding half the scale, then taking floor, rounds to the nearest integer with
exact half ties toward positive infinity. Numerator 16384 becomes one, while
$-16384$ becomes zero. Numerator $-49152$ represents $-1.5$ and becomes $-1$.
A division that truncates toward zero cannot replace the floor operation for
negative biased numerators.

The function stores $R_i$ clamped to $[-32768,32767]$, and counts a clip only
when the rounded integer lies strictly outside that interval. At the positive
end, a value of 32767.5 rounds to 32768 and clips. At the negative end,
$-32768.5$ rounds to $-32768$ and does not clip. Both a true negative clip and
that exact half case store bytes \texttt{00 80}; the output bytes alone cannot
explain the clipping count.

Do not round the products separately. Eight products whose numerators are each
16384 have an exact sum of 131072. One final rounding returns four. Rounding
each half first would return eight. Nor may partial sums saturate before the
remaining products arrive: a later negative product can cancel an earlier
positive one. The preceding chapter's failure of associative saturated addition
applies here too.

\begin{samepage}
One product fits in a signed word. The largest is
$(-32768)(-32768)=1073741824=2^{30}$. Two such products sum to $2^{31}$, which
does not fit. Eight sum to $2^{33}=8589934592$. The most negative individual
product is $(-32768)(32767)=-1073709056$, so the eight-product sum lies between
$-8589672448$ and 8589934592. Adding the rounding bias also fits comfortably in
a signed doubleword. These bounds justify exact 64-bit accumulation, not exact
32-bit accumulation followed by a wider copy.
\par
\end{samepage}

In the scalar path, signed sample and coefficient values are reconstructed from
two bytes with \texttt{seh}, as in the mixer. Each \texttt{dmul} produces an
exact doubleword product under these bounds; \texttt{daddu} adds it to a
zero-initialized doubleword sum. The inner loop consumes exactly eight taps.
The outer loop then applies the shared rounding and clamp rule and moves to the
next window. The complete source retains this scalar path as \texttt{fir\_scalar}.

\section{Widen before a vector dot product can overflow}

MSA's \texttt{dotp\_s.w} multiplies corresponding signed halfwords and adds
neighboring products into word lanes. That sounds suitable for sixteen-bit
samples, but a pair can already exceed the signed word range. An all-minimum
pair produces word bit pattern \texttt{80000000}; copying that lane as a signed
word yields a negative value. The lost sign cannot be repaired by copying it
into a 64-bit general register afterward.

Instead, sign-extend each input and coefficient to a word \emph{before}
multiplication. Then \texttt{dotp\_s.d} combines pairs of signed word products
into doubleword lanes. The original values are still only sixteen-bit numbers;
the wider container gives their pair sums enough range.

The following coefficient preparation runs once for a nonempty call. Register
\texttt{\$a1} points to exactly sixteen readable coefficient bytes.
\Needspace{10\baselineskip}
\MARSCodeLabel{Prepare signed word coefficients}
\begin{minted}{mars-idea}
ld.b $w1, 0($a1)
clti_s.h $w2, $w1, 0
ilvr.h $w3, $w2, $w1
ilvl.h $w4, $w2, $w1
\end{minted}

\texttt{clti\_s.h} produces \texttt{ffff} for each negative halfword and zero
otherwise. Interleaving a sample halfword with that sign halfword builds a signed
word: $-32768$ becomes \texttt{ffff8000}, and 4096 becomes \texttt{00001000}.
For this little-endian lane arrangement, \texttt{ilvr.h} takes the low four
halfwords, interleaving the last source first; \texttt{ilvl.h} takes the high
four. Read the resulting registers as words, not as eight independent samples.
The values in \texttt{\$w3} and \texttt{\$w4} remain available for every window.

The per-output calculation performs the same widening for the current eight
samples. Only its input address changes between windows.
\Needspace{18\baselineskip}
\MARSCodeLabel{Form an exact eight-product numerator}
\begin{minted}{mars-idea}
ld.b $w0, 0($a0)
clti_s.h $w2, $w0, 0
ilvr.h $w5, $w2, $w0
ilvl.h $w6, $w2, $w0
dotp_s.d $w7, $w5, $w3
dotp_s.d $w8, $w6, $w4
addv.d $w7, $w7, $w8
copy_s.d $t0, $w7[0]
copy_s.d $t1, $w7[1]
daddu $t0, $t0, $t1
\end{minted}

Write $p_j=x_{i+j}c_j$. The first dot product gives doublewords
$[p_0+p_1,p_2+p_3]$; the second gives $[p_4+p_5,p_6+p_7]$. Their vector sum
contains $[p_0+p_1+p_4+p_5,p_2+p_3+p_6+p_7]$. The two copies and final
\texttt{daddu} finish $S_i$. This grouping is safe because all these intermediate
sums fit in signed doublewords and none has yet been rounded or saturated.

For eight minimum inputs and coefficients, each dot-product lane is 2147483648,
hexadecimal \texttt{0000000080000000}. Each combined lane is 4294967296, and
the final scalar sum is 8589934592. This example distinguishes widening before
arithmetic from widening an already damaged result.

\section{Quantize one output and move one window}

The common quantizer begins with the exact numerator in \texttt{\$t0}:
\Needspace{8\baselineskip}
\MARSCodeLabel{Round a wide numerator once}
\begin{minted}{mars-idea}
daddiu $t0, $t0, 16384
dsra $t0, $t0, 15
\end{minted}
The arithmetic shift implements floor division by $2^{15}$. The complete source
then compares this wide rounded result strictly above 32767 and below $-32768$,
selects an endpoint only when necessary, increments \texttt{\$v1} for that event,
and stores the low two bytes. The bias and shift must operate on the full
numerator; narrowing it before conversion would undo the earlier width proof.
The rounded result itself lies between $-262136$ and 262144, so it fits in a
signed word, but still needs the separate halfword-range comparisons.

The interface takes full input, coefficient and output pointers in
\texttt{\$a0/\$a1/\$a2}, output count zero through 64 in \texttt{\$a3}, and
returns original count and clips in \texttt{\$v0/\$v1}. Samples and coefficients
may begin at any byte address. The three selected slices are disjoint. Their
lengths are $2(n+7)$, 16 and $2n$ bytes. Saved registers, frame pointer and stack
pointer retain all their bits; other caller data remains unchanged. The routine
returns through the supplied continuation and performs no console I/O.

After one output, advance the input pointer by two bytes, advance output by two,
and decrement the remaining output count. The window moves by one sample even
though it contains eight. After $k$ outputs, the first $k$ destination samples
are complete, the clipping count describes exactly that prefix, and the input
pointer identifies original sample $k$. On the final iteration, the window ends
at original index $n+6$, the last selected input sample. No remainder path is
needed: every output always has an eight-sample window, even when the output
count is one. Count zero must return before reading even the coefficients.

Disjoint output is part of this contract. A store into future input could change
a later overlapping window. The mixer's exact-alias proof does not transfer to
this filter. Extra allocated padding also does not enlarge the permitted input
slice. The access argument follows the selected window and count, not whether an
address happens to be mapped.

\section{Continue a stream with explicit history}

There is no hidden state between calls. To split after $m$ outputs, start the
second call at original input sample $m$, a byte displacement of $2m$. That
second slice begins with the seven samples preceding its first new sample;
its first window is exactly the original window for output $m$. Pass the same
coefficients, an output pointer advanced by $2m$, and count $n-m$. The two output
sequences must concatenate to the whole-call result. Add their clipping counts
to recover the whole count.

For the sixteen-output impulse example, split after five. The second window
starts at original sample five, so it still contains the impulse at original
sample seven. Its first three outputs are 2000, then zero. Supplying seven new
zeros as its history would erase those contributions. A correct kernel cannot
recover history that its caller discarded.

This explicit-history interface is convenient for checking. A file-processing
caller can retain seven trailing samples and prepend them to the next block.
It must preserve chronological order and distinguish short final input from
zero-length output. An audio device, file parser and real-time deadline are
outside the numerical kernel; later chapters supply a file pipeline.

\section{Filter fewer positions independently}

Implement \texttt{decimate\_q15} in \texttt{fir-decimate.asm}. Use the same eight
weights, exact accumulation, rounding and clipping policy, but compute windows
starting at original indices $0,2,4,\ldots,2(n-1)$. Return $n$ output samples
and their strict clipping count. The output count is zero through 32. The input
has seven history samples and, for a nonempty call, $2n-1$ new samples: $2n+6$
samples in total. A zero-output call declares seven history samples and reads
none of them. Coefficients occupy sixteen bytes; output occupies $2n$ bytes.
The pointers, disjointness, preservation and return rules match the worked filter.

This operation filters and retains every second position, called decimation by
two. It is not equivalent to discarding every other input before filtering;
that would change the values within each window. For the default averaging
impulse and eight requested outputs, the first four outputs are 2000 and the
remaining four are zero. With asymmetric weights, derive the even-indexed
outputs from the full filter instead of assuming another four identical values.

Prove the last window's input bound and derive the input-byte advance per output.
Then decide what input base a second call needs after $m$ retained outputs.
Establish a scalar baseline, check an equivalent vector implementation, and
investigate a negative half tie and an overflowing word pair separately. The
written lab supplies public checks; it does not require particular temporary
registers or instruction spelling. If vectorization is assigned, explain the
lane meanings and verify them in the debugger in addition to checking results.

\section*{Sources and laboratory connection}
The public \path{packages/021-fixed-point-fir/} supplies the original exact
contract, split-call checks and optional generated listening comparisons.
The MSA widening and interleave behavior follows the pinned MARS implementations.
Arm's \href{https://arm-software.github.io/CMSIS-DSP/latest/group__FIR.html}{CMSIS-DSP FIR documentation}
provides a separate example of Q15 filtering with a 64-bit accumulator. Its
conversion policy is not this exercise's nearest-rounding contract. Keep such
policies explicit when comparing implementations. The revised lab follows the
wide intermediate states and assigns the decimating filter independently.
